%% ---------------------------------------------------------------------------
%% DRAFT Methods section (referee item M7). Not yet \input by manuscript.tex.
%% Every measured number is a macro from paper/numbers.tex. Design constants are
%% literals, each with a `% source:` line naming the file and function it is read
%% from. \PENDINGMACRO marks a measured number that has no macro yet; it must be
%% added to paper/numbers.py before this section is merged.
%% ---------------------------------------------------------------------------
\providecommand{\PENDINGMACRO}[1]{\textbf{[PENDING: #1]}}

\section{Methods}
\label{sec:methods}

\subsection*{Governing model}

% source: solver_fd.py:generate_breakthrough_data (rhs, laplacian, _upwind_div)
The reference is a one-dimensional, non-isothermal packed bed with axial dispersion,
linear-driving-force (LDF) kinetics and a lumped energy balance with wall loss. With
\(c(z,t)\) the gas-phase adsorbate concentration, \(q(z,t)\) the adsorbed loading per
unit particle mass, \(T(z,t)\) the bed temperature and \(v\) the \emph{superficial}
velocity, the equations integrated are
% source: solver_fd.py:generate_breakthrough_data -- dc_dt = D_L*lap - adv_c(u_gas=v/eps_t) - src_factor*dq_dt
% TODO-VERIFY: the solver_fd.py module docstring writes the mass balance as
%   eps_t dc/dt = D_L d2c/dz2 - v dc/dz - (1-eps_t) rho_p dq/dt,
% which differs from the code by a factor eps_t on the dispersion term. The equation
% below is what rhs() actually integrates; the docstring should be corrected.
\begin{align}
\frac{\partial c}{\partial t} &= D_L \frac{\partial^2 c}{\partial z^2}
 - \frac{v}{\varepsilon_t}\frac{\partial c}{\partial z}
 - \frac{1-\varepsilon_t}{\varepsilon_t}\,\rho_p \frac{\partial q}{\partial t},
\label{eq:mass}\\
% source: solver_fd.py:generate_breakthrough_data -- dq_dt = k_LDF*(qstar_fn(c,T) - q)
\frac{\partial q}{\partial t} &= k_{\mathrm{LDF}}\bigl(q^{*}(c,T) - q\bigr),
\label{eq:ldf}\\
% source: solver_fd.py:generate_breakthrough_data -- alpha_T, u_thermal, heat_factor, wall_factor
C_{\mathrm{b}}\frac{\partial T}{\partial t} &= k_z \frac{\partial^2 T}{\partial z^2}
 - v\rho_g C_{pg}\frac{\partial T}{\partial z}
 + (1-\varepsilon_t)\rho_p(-\Delta H)\frac{\partial q}{\partial t}
 - \frac{4h_w}{D_{\mathrm{in}}}\,(T - T_w),
\label{eq:energy}
\end{align}
% source: solver_fd.py:AdsorptionPhysicsConfig.C_term
with bed heat capacity \(C_{\mathrm{b}} = \varepsilon_t\rho_g C_{pg} + (1-\varepsilon_t)\rho_p C_{ps}\)
and \(q^{*}\) the dual-site isotherm of equation~\eqref{eq:isotherm}. The cluster-to-primary
affinity ratio is fixed, \(b_{C0}/b_{H0} = 20\), and \(b_{C0}\) is set per run so that
the cooperative step sits at the material's step humidity at the feed temperature,
\(b_C(T_{\mathrm{in}})\,c_{\mathrm{step}} = 1\).
% source: gen_parametric_dataset.py:B_H_RATIO (=20.0) and build_physics (p.b_H0 = p.b0 / B_H_RATIO)
% source: isotherm.py:calibrate_step; gen_parametric_dataset.py:build_physics (c_step = rh_to_conc(step_rh, T_in))
Feed concentration follows from relative humidity through the ideal-gas law and a
Magnus saturation pressure, \(c_{\mathrm{in}} = \mathrm{RH}\,p_{\mathrm{sat}}(T_{\mathrm{in}})/(RT_{\mathrm{in}})\).
% source: gen_parametric_dataset.py:rh_to_conc, p_sat_water (610.94*exp(17.625*Tc/(Tc+243.04)))

\paragraph{Boundary and initial conditions.} The inlet is of flux (Danckwerts) type for
both mass and heat: the feed enters advectively and the inlet face carries no
dispersive or conductive flux,
% source: solver_fd.py:generate_breakthrough_data.laplacian (lap[0] = (f[1]-f[0])/dz^2) and _upwind_div (f_in at the inlet face)
\begin{equation}
\frac{v}{\varepsilon_t}\,c_{\mathrm{in}} = \frac{v}{\varepsilon_t}\,c - D_L\frac{\partial c}{\partial z},
\qquad
v\rho_g C_{pg}\,T_{\mathrm{in}} = v\rho_g C_{pg}\,T - k_z\frac{\partial T}{\partial z}
\qquad (z = 0),
\label{eq:danckwerts}
\end{equation}
% source: solver_fd.py:generate_breakthrough_data.laplacian (lap[-1] = (f[-2]-f[-1])/dz^2)
with zero gradient, \(\partial c/\partial z = \partial T/\partial z = 0\), at the outlet
\(z = L\). The bed starts clean and at wall temperature, \(c = q = 0\) and \(T = T_w\),
and the wall is held at the feed temperature, \(T_w = T_{\mathrm{in}}\).
% source: solver_fd.py:generate_breakthrough_data (y0 = [0, 0, T_w]); gen_parametric_dataset.py:build_physics (p.T_w = cond["T_in"])

\paragraph{Fixed bed and transport properties.}
% source: gen_parametric_dataset.py:build_physics (p.L = 0.10; p.D_L = 0.7*2.5e-5 + 0.5*d_p*(v/eps_t))
% source: solver_fd.py:AdsorptionPhysicsConfig.__init__ (D_in, rho_g, C_pg, C_ps, k_z, h_w)
Every run uses \(L = 0.10\)\,m, \(D_{\mathrm{in}} = 0.05\)\,m,
\(\rho_g = 1.2\)\,kg\,m\(^{-3}\), \(C_{pg} = C_{ps} = 1000\)\,J\,kg\(^{-1}\)\,K\(^{-1}\),
\(k_z = 0.10\)\,W\,m\(^{-1}\)\,K\(^{-1}\) and \(h_w = 10\)\,W\,m\(^{-2}\)\,K\(^{-1}\). Axial
dispersion follows the packed-bed correlation \(D_L = 0.7D_m + 0.5\,d_p\,v/\varepsilon_t\)
with \(D_m = 2.5\times10^{-5}\)\,m\(^2\)\,s\(^{-1}\).
% TODO-VERIFY: the dispersion correlation uses D_m = 2.5e-5 (gen_parametric_dataset.py:build_physics)
% while the Glueckauf relation uses D_M_WATER = 2.6e-5 (gen_parametric_dataset.py:D_M_WATER).
% Small, but two values of one constant in one generator.
In design v2 the kinetic coefficient is not sampled; it is derived from the particle
through the Glueckauf relation using the chord slope of the isotherm at the feed,
% source: gen_parametric_dataset.py:k_ldf_glueckauf; constants EPS_PELLET=0.35, TORTUOSITY=3.0, D_M_WATER=2.6e-5
\begin{equation}
k_{\mathrm{LDF}} = \frac{15\,\varepsilon_p D_m/\tau_p}{(d_p/2)^2\,\rho_p\,K},
\qquad K = \frac{q^{*}(c_{\mathrm{in}},T_{\mathrm{in}})}{c_{\mathrm{in}}},
\label{eq:glueckauf}
\end{equation}
with pellet porosity \(\varepsilon_p = 0.35\), tortuosity \(\tau_p = 3\) and
\(D_m = 2.6\times10^{-5}\)\,m\(^2\)\,s\(^{-1}\). One \(k_{\mathrm{LDF}}\) is therefore
constant within a run but differs between conditions of the same material.

\paragraph{Numerics and storage.}
% source: gen_parametric_dataset.py:SOLVE_NZ (2000), run_one (n_snapshots=600); solver_fd.py:generate_breakthrough_data (BDF, rtol=1e-6, atol=1e-10, jac_sparsity)
Equations~\eqref{eq:mass}--\eqref{eq:energy} are discretised by the method of lines on
2000 cell-centred finite volumes (first-order upwind advection, second-order central
dispersion and conduction) and integrated with SciPy's BDF method at relative and
absolute tolerances \(10^{-6}\) and \(10^{-10}\), with an explicit Jacobian sparsity
pattern, to 600 output times. The verification of this solver is in
Section~\ref{sec:reference}.
% source: gen_parametric_dataset.py:HORIZON (2.5) and run_one (multipliers HORIZON, 2*, 5*, 12*; exit_ratio >= 0.90)
Each run is integrated to \(t_{\mathrm{final}} = 2.5\,t_{\mathrm{st}}\), with
\(t_{\mathrm{st}} = [(1-\varepsilon_t)\rho_p q^{*}(c_{\mathrm{in}},T_{\mathrm{in}}) + \varepsilon_t c_{\mathrm{in}}]\,L/(v c_{\mathrm{in}})\)
the stoichiometric time; if the outlet has not reached \(0.90\,c_{\mathrm{in}}\) the
horizon is extended to 5, 12.5 and 30 \(t_{\mathrm{st}}\) in turn and the run is
rejected only if all four fail.
% source: solver_fd.py:AdsorptionPhysicsConfig.stoichiometric_time
% source: gen_parametric_dataset.py:screen (theta in [0.10, 0.995]; K_H finite and > 0; MTZ >= 5 cells; 60 s < t_st < 2e6 s)
Before solving, a draw is rejected if the feed loads less than 0.10 or more than 0.995
of capacity, if the Henry constant is not finite and positive, if the isothermal
mass-transfer zone spans fewer than five solver cells, or if \(t_{\mathrm{st}}\) lies
outside 60\,s to \(2\times10^{6}\)\,s; \nDataNrejected\ of the design-v2 draws were
rejected this way and \nDataNsamples\ accepted.
% source: gen_parametric_dataset.py:run_one (fields = c/c_in, q/q_max, T/T_in, bilinear resample); chain_gen_v2.sh (ADS_STORE_NZ=256, ADS_STORE_NT=256)
Fields are stored as \(c/c_{\mathrm{in}}\), \(q/q_{\max}\) and \(T/T_{\mathrm{in}}\) on a
uniform grid in \(z/L\) and \(t/t_{\mathrm{final}}\), bilinearly resampled to
\(256\times256\) for design v2 and \(512\times256\) for the legacy design, with
\(t_{\mathrm{final}}\) kept as a per-sample scalar.

\subsection*{Damk\"ohler number}

% source: gen_parametric_dataset.py:run_one ("Da": p.k_LDF * t_final); dataset_summary.py:summarise; run_l6_v2.py:damkohler_per_material
The Damk\"ohler number used throughout is computed per sample as
\begin{equation}
\mathrm{Da} = k_{\mathrm{LDF}}\; t_{\mathrm{final}},
\label{eq:da}
\end{equation}
the LDF rate multiplied by the simulated horizon. The per-material value on which rung
L6 regresses is the median of Equation~\eqref{eq:da} over that material's accepted
conditions. Because \(t_{\mathrm{final}}\) is a multiple of \(t_{\mathrm{st}}\) that
depends on whether the horizon was extended, Equation~\eqref{eq:da} is \(2.5\,k_{\mathrm{LDF}}t_{\mathrm{st}}\)
for unextended runs and a larger multiple for the
\PENDINGMACRO{fraction of design-v2 samples whose horizon was extended beyond 2.5 t_st, data/parametric_v2/manifest.json, samples[].horizon_extensions > 0}
of samples that were extended.
% TODO-VERIFY: Da therefore carries a factor of the horizon multiplier (2.5, 5 or 12.5;
% counted from the manifest at drafting time, 0/1/2 extensions) in addition to the
% kinetic ratio k_LDF*t_st. The paper treats Da as a kinetic-regime axis; decide
% whether to keep this definition (and say so) or switch to k_LDF*t_st.
% TODO-VERIFY: pde_adsorption.py:NondimConfig.groups and run_l4.py:build_phys_table
% define a DIFFERENT "Da" = k_LDF * L / v (convective time). That group is internal
% to the PINN residual and never reported, but it shares the name.

\subsection*{Parameter space}

% source: gen_parametric_dataset.py:MATERIAL_SPACE_V2, CONDITION_SPACE, sample_design, _sobol_unit; chain_gen_v2.sh (--design v2 --materials 240 --conditions 17 --seed 20260830)
Materials and operating conditions are sampled separately and crossed, 240 materials
at 17 conditions each, so that a material is a real grouping in the data
(Table~\ref{tab:params}). Each set is a scrambled Sobol design (SciPy \texttt{qmc.Sobol},
the next power of two truncated to the required count) mapped \emph{linearly} onto the
ranges; no parameter is sampled on a logarithmic scale. Materials use seed 20260830 and
conditions seed 20260837.
% source: gen_parametric_dataset.py:main (conditions seed = args.seed + 7)
The legacy design used by rungs L3 and L5 and the coordinate change differs in three
ways: independent uniform draws, \(k_{\mathrm{LDF}}\) sampled directly on
\([0.002, 0.05]\)\,s\(^{-1}\) in place of \(d_p\), and \(d_p\) fixed at 2\,mm; it has
60 materials at 17 conditions and \PENDINGMACRO{accepted legacy simulations, data/parametric/manifest.json, sum of counts}
accepted runs.
% source: gen_parametric_dataset.py:MATERIAL_SPACE (k_LDF 0.002-0.05), sample_material (rng.uniform), build_physics (d_p default 0.002); data/parametric/manifest.json (60 materials, 17 conditions, seed 20260820)

\begin{table}[t]
\centering
\caption{The design-v2 parameter space. All eleven are sampled by scrambled Sobol and
mapped linearly. Every design-v2 surrogate except rung L6 receives exactly these eleven
values, standardised on its own fitting set; L6 receives observable descriptors instead.}
\label{tab:params}
% source: gen_parametric_dataset.py:MATERIAL_SPACE_V2 and CONDITION_SPACE; ladder_data.py:PARAM_KEYS_V2 (input order)
\begin{tabular}{llll}
\toprule
parameter & symbol & units & range \\
\midrule
\multicolumn{4}{l}{\emph{material} (8)} \\
saturation capacity & \(q_{\max}\) & mol\,kg\(^{-1}\) & 8--30 \\
heat of adsorption & \(\Delta H\) & kJ\,mol\(^{-1}\) & \(-60\) to \(-38\) \\
step position (relative humidity) & \(\mathrm{RH}_{\mathrm{step}}\) & -- & 0.08--0.45 \\
cooperativity exponent & \(n\) & -- & 1--6 \\
Henry-site fraction & \(f_H\) & -- & 0.03--0.25 \\
particle diameter & \(d_p\) & mm & 1.5--5.0 \\
particle density & \(\rho_p\) & kg\,m\(^{-3}\) & 700--1400 \\
bed porosity & \(\varepsilon_t\) & -- & 0.30--0.48 \\
\midrule
\multicolumn{4}{l}{\emph{operating condition} (3)} \\
feed relative humidity & \(\mathrm{RH}_{\mathrm{feed}}\) & -- & 0.15--0.85 \\
superficial velocity & \(v\) & m\,s\(^{-1}\) & 0.03--0.30 \\
feed temperature & \(T_{\mathrm{in}}\) & K & 288.15--313.15 \\
\bottomrule
\end{tabular}
\end{table}

\subsection*{Error metric and statistical test}

% source: metrics.py:per_variable_nrmse; v2_common.py:sample_row
The primary metric is a per-sample normalised root-mean-square error on the
\emph{concentration} channel. For sample \(s\), with \(\hat c_{ij}\) and \(c_{ij}\) the
predicted and reference \(c/c_{\mathrm{in}}\) on the rung's evaluation grid
(\(i\) over \(N_z\) positions, \(j\) over \(N_t\) times),
\begin{equation}
\mathrm{nRMSE}_c(s) = \frac{\Bigl[\frac{1}{N_zN_t}\sum_{i,j}(\hat c_{ij}-c_{ij})^2\Bigr]^{1/2}}
{\max_{i,j} c_{ij} - \min_{i,j} c_{ij}} ,
\label{eq:nrmse}
\end{equation}
the RMSE over the whole space--time field divided by that sample's own reference range.
% source: gen_parametric_dataset.py:run_one (initial clean bed; exit_ratio >= 0.90 acceptance)
Because every accepted run starts from a clean bed and ends with the outlet above
\(0.90\,c_{\mathrm{in}}\), the denominator is at least 0.90, so
Equation~\eqref{eq:nrmse} is in effect an RMSE in units of the feed concentration and
cannot be inflated by a vanishing range.
% source: run_l4b.py:eval_window (range over the FULL reference field, not the window -- A15)
Where only part of the field is scored (the time-extrapolation axis of rung L4), the
numerator runs over the scored window but the denominator remains the range of the full
reference field.
% source: metrics.py:pooled_mse (raises), per_variable_nrmse (q, T channels)
The same quantity is computed separately for \(q/q_{\max}\) and \(T/T_{\mathrm{in}}\) and
reported as secondary; the three channels are never pooled, because the temperature
variance would dominate.
% source: v2_common.py:exit_nrmse (range < 0.05 raises); run_l7_v2.py (exit curve, 128 times)
Rung L7 scores the outlet curve only: Equation~\eqref{eq:nrmse} restricted to \(z = L\),
normalised by the range of the reference outlet curve, which is required to exceed 0.05.
% source: analyze_l1_v2.py:seed_avg, pooled_folds
Per-sample errors are averaged over the three training seeds, pooled across held-out
folds so that each sample appears once, and averaged with equal weight per sample; a
material therefore contributes in proportion to its number of accepted conditions.

% source: metrics.py:_cluster_bootstrap, compare; analyze_*.py (n_boot=4000)
Two arms are compared on identical samples by a paired, cluster-robust percentile
bootstrap: materials are resampled with replacement, all conditions of each drawn
material are kept, and the mean paired difference is recomputed 4000 times. A
difference is reported as significant only if the \((1-\alpha)\) percentile interval
excludes zero. The level \(\alpha\) is calibrated per design rather than set at 0.05.
% source: calibrate_v2.py:simulate (400 trials, n_boot 600, real cluster sizes) and main (grid 0.05, 0.02, 0.01, 0.005, 0.002; first with FPR <= 5 %)
Using each design's real conditions-per-material counts, we simulate clustered nulls in
which the arms tie on average but individual materials favour one or the other, and take
the largest nominal level in \(\{0.05, 0.02, 0.01, 0.005, 0.002\}\) whose empirical
false-positive rate does not exceed 5\,\%. This gives \(\alpha = \nAlphaFolds\) for the
\nNclustFolds-material five-fold design (empirical size \nFprFoldsCal\,\%) and
\(\alpha = \nAlphaManifest\) for the \nNclustManifest-material split; on the
\nNclustLegacy-material legacy design we report at the stricter \(\alpha = \nLfiveAlpha\)
(empirical size \nFprLegacyUsed\,\%).
% source: run_l6_v2.py:slope_test (material bootstrap of the OLS slope, 4000 draws, ALPHA_V2=0.02)
Rung L6's primary estimand is instead the ordinary-least-squares slope of the
per-material paired difference on \(\log_{10}\mathrm{Da}\), with a material bootstrap
interval at \(\alpha = \nAlphaFolds\).
% TODO-VERIFY: L4's time axis has 192 clusters and is tested at alpha = 0.02
% (analyze_l4b_v2.py:ALPHA), a level calibrated at 48 and 240 clusters but not at 192.
% Both neighbouring calibrations select 0.02, so this is likely harmless; say so or calibrate.

\subsection*{Surrogates, encodings and evaluation designs}

% source: v2_common.py:load_folds; results/l6_v2_folds.json (n_folds 5, fold_seed 20260901)
Three evaluation designs are used, and each result states which. \emph{Five-fold}: the
240 design-v2 materials are partitioned into five folds (seed 20260901), every material
is held out once, and each fit sees 192 materials. \emph{Single split}: the dataset's own
manifest split, holding out \nNclustManifest\ materials with 192 for training. \emph{Legacy}: the earlier dataset's split, holding out
\nNclustLegacy\ of 60 materials.
% source: gen_parametric_dataset.py:main (holdout-materials 0.20, rng seed+1); ladder_data.py:load
Every learned arm is trained with seeds 42, 43 and 44, sees the same eleven
standardised parameters (except L6) and is fitted, including any proper orthogonal
decomposition (POD) basis and scaler, on the fit set only.
% source: v2_common.py:assert_disjoint, standardise_params; run_l1.py:fit_pod
Table~\ref{tab:rungs} gives the architecture and training of each rung.

% source: run_l5.py module docstring ("Cross-rung comparison against L1-L4 absolute numbers is NOT valid and is not made.")
\textbf{Absolute errors are not comparable across rungs.} The rungs differ in dataset
(design v2 or legacy), in field encoding (\(128\times128\) subsampled from
\(256\times256\), \(128\times128\) subsampled from \(512\times256\), or the full
\(512\times256\)), in what is scored (the whole field, a time window, or the outlet
curve) and in held-out set. An L3 error of \nLthreeMlpLarge\ and an L5 error of
\nLfiveONet\ are measured on the same legacy materials but on different encodings ---
L3 on the full \(512\times256\) field, L5 on a \(128\times128\) subsample of it --- and
their ratio is not an architecture comparison. Only paired comparisons within a rung are
inferential; Table~\ref{tab:rungs}
states the encoding of each so that no cross-rung reading is made by accident.
% source: ladder_data.py:load (field_res index subsampling); run_l3.py:main (ladder_data.load() at stored 512x256); run_l5.py:subsample (NZ_S=NT_S=128); comoving.py:NZ_S, load_c_light

\begin{table}[p]
\centering
\caption{Surrogates by rung. ``Five-fold'', ``single split'' and ``legacy'' are the
three evaluation designs of the text. Field sizes are \(N_z\times N_t\) on the
evaluation grid. Parameter counts and widths are design choices fixed before training.
Absolute errors in different rows are not comparable.}
\label{tab:rungs}
\footnotesize
\renewcommand{\arraystretch}{1.15}
\begin{tabular}{@{}p{0.07\textwidth}p{0.13\textwidth}p{0.24\textwidth}p{0.29\textwidth}p{0.19\textwidth}@{}}
\toprule
rung & data, design & input \(\to\) output; encoding & architecture and size & training \\
\midrule
% source: run_l1_v2.py:main (--modes 64, --seeds 42 43 44, FIELD_RES 128); run_l1.py:make_arm, fit_pod
L1 & v2, five-fold &
11 parameters \(\to\) POD coefficients, 64 modes per channel for \(c,q,T\); field \(128\times128\) &
ridge (\(\lambda=1\)); random forest (300 trees, leaf 2); gradient boosting (400 trees, depth 5, rate 0.06); perceptron \(3\times256\), ReLU, standardised targets &
scikit-learn; perceptron Adam \(10^{-3}\), early stopping (patience 30, at most 3000 epochs) \\
% source: learning_curve_v2.py:SIZES (12, 24, 48, 96, 192)
L1 curve & v2, five-fold & as L1, training materials subsampled to 12, 24, 48, 96, 192 & the L1 perceptron & as L1 \\
% source: run_l2.py:sweep_configs; run_l2_v2.py:main
L2 & v2, five-fold & as L1 &
28 configurations: perceptron width 16--1024 at depth 3; depth 1--10 at width 256; boosting depth 2--10; forest leaf 1--32 & as L1 \\
% source: run_l3.py:build, RBFKANBlock (grids=8), ChebyKANBlock (degree=7), solve_width, train_one; results/l3_results.json (steps 12000; widths per budget); results/l3_merged.json (grid 3e-5..1e-1)
L3 & legacy &
11 parameters (legacy: \(k_{\mathrm{LDF}}\) in place of \(d_p\)) \(\to\) POD coefficients, 64 per channel; full \(512\times256\) field &
depth 3; GELU perceptron, Gaussian-RBF KAN (8 centres per edge), Chebyshev KAN (degree 7); widths solved to 50k / 200k / 800k parameters within \(\pm10\,\%\): perceptron 114/269/583, RBF 45/98/203, Chebyshev 48/103/215 &
Adam, cosine to \(\mathrm{lr}/100\), clip 1, 12\,000 full-batch steps, best of a 12\,\% validation split (patience 2000 steps); learning rate per family over \(3\times10^{-5}\)--\(10^{-1}\) \\
% source: run_l4b_v2.py:ARMS, main (width 192, depth 4, steps 8000, lr 1e-3, n_sup 4096, n_colloc 1024, n_bc 512, polish_iters 500); run_l4.py:ParametricPINN; results/l4b_v2_results.json (n_params 114435, n_eval_time 768, n_eval_material 793)
L4 & v2; time axis: 192 training materials, \(t/t_{\mathrm{final}}\in[0,0.5]\) seen, \([0.5,1]\) scored; material axis: single split &
(11 parameters, \(z\), \(t\)) \(\to\) \((c,q,T)\) pointwise, bounded output head; scored on \(128\times128\) &
GELU perceptron \(4\times192\) (114\,435 parameters); 13 arms: data-only twin, fixed residual weights \(10^{-6}\)--\(1\), gradient-norm balancing, NTK and self-adaptive weights; a nondimensional interior residual plus boundary and initial terms (see TODO-VERIFY below the table) &
Adam \(10^{-3}\), cosine, 8000 steps, 4096 data, 1024 collocation and 512 boundary points per step; then 500 L-BFGS iterations (polish) \\
% source: run_l5.py:DeepONet, RBFKANStack (grids 4), solve_width, subsample, train_eval; run_l5_fno.py:FNOArm, main (layers 4, modes 4/8/16, batch 8); results/l5_merged.json (budget 200000, steps 8000)
L5 & legacy &
11 parameters \(\to\) \(c\) field (all three channels trained, \(c\) scored); \(128\times128\) subsampled from \(512\times256\) &
DeepONet and DeepOKAN, branch and trunk depth 3, \(p\in\{8,\dots,128\}\); FNO, 4 Fourier layers, 4--16 modes; widths solved to 200k parameters (best DeepONet \nLfiveONetWidth, best FNO \nLfiveFNOwidth) &
Adam, cosine, clip 1, 8000 steps, 32 samples \(\times\) 2048 points (FNO: 8 full fields); \nLfiveNrates\ learning rates over \nLfiveDecades\ decades \\
% source: l5_bottleneck_v2.py (HistGradientBoostingRegressor max_iter 400, lr 0.06; channel c; 128x128; five folds)
L5 basis vs.\ map & v2, five-fold &
11 parameters \(\to\) POD coefficients of \(c\), \(p\in\{8,\dots,128\}\); \(128\times128\) &
one gradient-boosted regressor per mode (400 iterations, rate 0.06), standardised coefficients & three seeds \\
% source: run_l6.py:Joint, Separate, sample_rays, train; run_l6_v2.py:main (budget 120000, depth 4, steps 12000, lr 1e-3, n_rays 128, n_t 32); results/l6_v2_results.json (n_desc 54); descriptors.py:build
L6 & v2, five-fold &
54 observable descriptors (isotherm at three temperatures, bed and operating values; \(d_p\) and isotherm parameters withheld), \(z\), \(t\) \(\to\) \((c,q,T)\); \(128\times128\) &
joint: one perceptron; separate: equilibrium head \(q^{*}\), kinetic head \(\log k\), transport head, with \(q\) obtained by integrating Eq.~\eqref{eq:ldf}; both 120k parameters, depth 4 &
Adam \(10^{-3}\), cosine, 12\,000 steps, 128 rays \(\times\) 32 times per step \\
% source: run_l7_v2.py; classical.py:MODELS (equilibrium_shock, klinkenberg, constant_pattern)
L7 & v2, all samples; paired with L1 out of fold &
parameters \(\to\) outlet curve, 128 times &
equilibrium shock, Klinkenberg, constant pattern; closed form, nothing fitted & none \\
% source: warp_verdict.py:pod_hgb, predict_warp (LEVELS 0.05/0.95, P_SHAPE 32, P_WARP 4); comoving2.py:N_SIG 512, SIG_LO -1.5, SIG_HI 2.5; comoving.py:load_c_light
Warp & legacy &
11 parameters \(\to\) \(c\) only, \(128\times128\); fixed frame, or the frame \(\sigma=(t-t_{0.05}(z))/(t_{0.95}(z)-t_{0.05}(z))\) on 512 points of \([-1.5,2.5]\) &
32 POD modes, one gradient-boosted regressor per mode; front trajectories from 4 POD modes of (start, log span) &
three seeds; oracle arm given the true fronts \\
\bottomrule
\end{tabular}
\end{table}

% TODO-VERIFY (L4, important): the physics residual minimised by rung L4
% (run_l4.py:pde_residual, boundary_residual) does not match Eq. (mass)/(danckwerts)
% as solved by solver_fd.py. In the code's own scaling (Lambda, Pe = vL/D_L, tau =
% t_final v / L) the residual's advection coefficient is tau/Lambda, whereas Eq. (mass)
% nondimensionalised the same way gives tau/(eps_t*Lambda): the residual advects at the
% superficial speed v, the solver at the interstitial speed v/eps_t (a factor 2.1-3.3
% over the eps_t range). The inlet residual c* - (1/Pe) dc*/dz* = 1 likewise lacks the
% eps_t of the solver's Danckwerts condition (eps_t/Pe), and the thermal inlet is
% Dirichlet (T* = 1) in the residual but flux-type in the solver. Dispersion, kinetics
% and energy-interior terms match. If confirmed, the residual is not the PDE the data
% obey, which bears directly on the "physics at inference makes transfer worse" result.
% Until resolved, the L4 row above should not say "residual of Eqs. (mass)-(energy)".
% TODO-VERIFY: run_l5_v2.py (DeepONet re-run on the five-fold design) is in progress;
% no macro exists for it, and this section does not report it.
